\section*{§1. Первое начало термодинамики}
\addcontentsline{toc}{section}{§1. Первое начало термодинамики}

\begin{basics}
\sub{Уравнение состояния идеального газа (Клапейрон–Менделеев)}
\[pV=\nu RT=\frac m\mu RT=NkT,\qquad p=nkT.\]
$p$ — давление (Па), $V$ — объём (м$^3$), $T$ — абсолютная температура (К);
$\nu=m/\mu$ — число молей, $m$ — масса газа, $\mu$ — молярная масса (кг/моль);
$N$ — число молекул, $n=N/V$ — концентрация (м$^{-3}$);
$R=8{,}31$ Дж/(моль\,К), $k=R/N_A=1{,}38\cdot10^{-23}$ Дж/К, $N_A=6{,}02\cdot10^{23}$ моль$^{-1}$.
Смесь: $p=\sum p_i$ (закон Дальтона), $p_i$ — давление компоненты, если бы она была одна.

\sub{Внутренняя энергия идеального газа}
\[U=\frac i2\nu RT=\nu C_VT,\qquad C_V=\frac i2R.\]
$i$ — число степеней свободы молекулы: 3 (одноатомный: He, Ar), 5 (двухатомный: H$_2$, N$_2$, O$_2$, воздух), 6 (многоатомный).
$U$ зависит \emph{только от $T$} — поэтому $\Delta U=\nu C_V\Delta T$ в \emph{любом} процессе, не только изохорном.

\sub{Работа газа}
\[\delta A=p\,\dd V,\qquad A=\int_{V_1}^{V_2}p(V)\,\dd V\]
— площадь под графиком на $p$--$V$. При расширении $A>0$, при сжатии $A<0$. Зависит от пути.

\sub{Первое начало термодинамики}
\[\delta Q=\dd U+\delta A,\qquad Q=\Delta U+A.\]
$Q$ — теплота, \emph{полученная} газом ($Q<0$ — газ отдал тепло). Это закон сохранения энергии.
$\dd U$ — полный дифференциал (функция состояния), $\delta Q$ и $\delta A$ — нет (зависят от пути).
За цикл $\Delta U=0\Rightarrow A=\sum Q$.

\sub{Теплоёмкость}
$\tilde C=\delta Q/\dd T$ — теплоёмкость тела; молярная $C=\tilde C/\nu$, удельная $c=\tilde C/m$.
Зависит от процесса! Для моля идеального газа из первого начала:
\[C=\frac{\delta Q}{\dd T}=C_V+p\frac{\dd V}{\dd T}\qquad(\text{производная — вдоль процесса}).\]
Изобара: $p\frac{\dd V}{\dd T}=R\Rightarrow C_p=C_V+R$ (уравнение Майера).
Показатель адиабаты $\gamma=\dfrac{C_p}{C_V}=\dfrac{i+2}{i}$: $\frac53$, $\frac75$, $\frac43$; отсюда $C_V=\dfrac{R}{\gamma-1}$, $C_p=\dfrac{\gamma R}{\gamma-1}$.

\sub{Изопроцессы (на $\nu$ молей)}
\begin{tabular}{@{}llll@{}}
процесс & $A$ & $Q$ & $C$\\\hline
изохора $V$=const & $0$ & $\nu C_V\Delta T$ & $C_V$\\
изобара $p$=const & $p\Delta V=\nu R\Delta T$ & $\nu C_p\Delta T$ & $C_p$\\
изотерма $T$=const & $\nu RT\ln\frac{V_2}{V_1}$ & $Q=A$ ($\Delta U=0$) & $\infty$\\
адиабата $Q=0$ & $-\Delta U=\nu C_V(T_1-T_2)$ & $0$ & $0$
\end{tabular}

\sub{Адиабата (обратимая)}
\[pV^\gamma=\text{const},\qquad TV^{\gamma-1}=\text{const},\qquad p^{1-\gamma}T^\gamma=\text{const},\]
\[A=\frac{p_1V_1-p_2V_2}{\gamma-1}=\nu C_V(T_1-T_2).\]
Получается из $C=0$: $C_V\dd T+p\,\dd V=0$. Адиабата на $p$--$V$ круче изотермы.

\sub{Политропа ($C=\text{const}$, обобщает все изопроцессы)}
\[pV^n=\text{const},\qquad n=\frac{C-C_p}{C-C_V},\qquad C=C_V\frac{n-\gamma}{n-1},\]
\[A=\frac{p_1V_1-p_2V_2}{n-1},\qquad Q=\nu C\Delta T.\]

\sub{Циклы}
$Q_1$ — тепло, полученное от нагревателя (сумма всех $Q_i>0$), $Q_2$ — отданное холодильнику ($Q_i<0$).
\[\eta=\frac{A}{Q_1}=1-\frac{|Q_2|}{Q_1},\qquad\text{Карно: }\eta=1-\frac{T_2}{T_1};\]
\[\text{холодильник: }\eta_x=\frac{|Q_2|}{A}\ (\text{может быть}>1).\]
Цикл по часовой стрелке на $p$--$V$ — тепловая машина ($A>0$), против — холодильная.
\end{basics}

%-------------------------------------------------------------
\begin{problem}{1.1}{exdz}{С одним молем идеального газа проводят процесс $p=p_0-\beta V^2$, где $p_0$ и $\beta>0$ — постоянные.
Найти максимально возможную температуру газа в этом процессе.}
\textbf{1. Выражаем $T$ через $V$.} Параметры $p,V,T$ связаны уравнением состояния $pV=RT$ и уравнением процесса, поэтому независим один — берём $V$:
\[T(V)=\frac{pV}{R}=\frac{(p_0-\beta V^2)V}{R}=\frac{p_0V-\beta V^3}{R}.\]

\textbf{2. Ищем максимум:} $\frac{\dd T}{\dd V}=\frac{p_0-3\beta V^2}{R}=0\Rightarrow V_m=\sqrt{\frac{p_0}{3\beta}}$.

\textbf{3. Подставляем:} $p(V_m)=p_0-\frac{p_0}3=\frac{2p_0}3$, $T_{\max}=\frac{p(V_m)V_m}{R}$.
\answer{$T_{\max}=\dfrac{2p_0}{3R}\sqrt{\dfrac{p_0}{3\beta}}$ (при $V=\sqrt{p_0/3\beta}$).}
\note{так же: 1.14 ($p=p_0e^{-\beta V}$, $T_{\max}=\frac{p_0}{\beta eR}$), 1.15 (минимум $p$).}
\end{problem}

%-------------------------------------------------------------
\begin{problem}{1.2}{ex}{Для идеального газа при $0^\circ$C найти термические коэффициенты: расширения $\alpha=\frac1{V_0}\big(\frac{\partial V}{\partial T}\big)_p$,
давления $\lambda=\frac1{p_0}\big(\frac{\partial p}{\partial T}\big)_V$ и сжатия $\beta=-\frac1{V_0}\big(\frac{\partial V}{\partial p}\big)_T$. Индекс 0 — при $T_0=273$\,К.}
Дифференцируем $pV=\nu RT$, держа постоянным указанный параметр:
\begin{itemize}
\item $p=$\,const: $p\,\dd V=\nu R\,\dd T\Rightarrow\big(\frac{\partial V}{\partial T}\big)_p=\frac{\nu R}{p}=\frac{V}{T}$, т.е.\ $\alpha=\frac1{T_0}$;
\item $V=$\,const: $\big(\frac{\partial p}{\partial T}\big)_V=\frac{\nu R}{V}=\frac pT$, т.е.\ $\lambda=\frac1{T_0}$;
\item $T=$\,const: $V=\frac{\nu RT}p\Rightarrow\big(\frac{\partial V}{\partial p}\big)_T=-\frac{V}{p}$, т.е.\ $\beta=\frac1{p_0}$.
\end{itemize}
\answer{$\alpha=\lambda=1/T_0$,\; $\beta=1/p_0$; выполняется $\alpha=p_0\beta\lambda$.}
\end{problem}

%-------------------------------------------------------------
\begin{problem}{1.3}{exdz}{Определить массу воздуха между двумя оконными рамами при атмосферном давлении $p_0=10^5$\,Па, если температура между рамами
меняется по линейному закону от $T_1=270$\,К до $T_2=300$\,К. Площадь окна $S=3$\,м$^2$, расстояние между рамами $l=0{,}2$\,м, $\mu=29$\,г/моль.}
\textbf{Идея:} газ неоднороден, но тонкий слой $\dd x$ можно считать однородным.

\textbf{1.} Ось $x$ перпендикулярно рамам: $T(x)=T_1+\frac{T_2-T_1}{l}x$. Для слоя толщиной $\dd x$ (объём $S\,\dd x$) уравнение состояния:
$p_0S\,\dd x=\frac{\dd m}{\mu}RT(x)\Rightarrow\dd m=\frac{\mu p_0S}{R}\cdot\frac{\dd x}{T(x)}$.

\textbf{2. Интегрируем} от 0 до $l$ (интеграл вида $\int\frac{\dd x}{a+bx}=\frac1b\ln(a+bx)$):
\[m=\frac{\mu p_0S}{R}\int_0^l\frac{\dd x}{T_1+(T_2-T_1)x/l}=\frac{\mu p_0Sl}{R(T_2-T_1)}\ln\frac{T_2}{T_1}.\]

\textbf{3. Числа:} $\frac{0{,}029\cdot10^5\cdot3\cdot0{,}2}{8{,}31\cdot30}=6{,}98$\,кг, $\ln\frac{300}{270}=0{,}105$.
\answer{$m=\dfrac{\mu p_0Sl}{R(T_2-T_1)}\ln\dfrac{T_2}{T_1}\approx0{,}735$ кг.}
\note{при $T_2\to T_1$ получаем обычное $m=\frac{\mu pV}{RT}$ — проверка.}
\end{problem}

%-------------------------------------------------------------
\begin{problem}{1.4}{ex}{Моль идеального газа ($i=5$) три раза обратимо переводят из состояния 1 $(p_1,V_1)$ в состояние 2 $(2p_1,2V_1)$ путями 1--3--2, 1--4--2 и 1--2
(рис.~1.5: 3 — $(2p_1,V_1)$, 4 — $(p_1,2V_1)$, 1--2 — отрезок прямой). Найти $Q_{132}$, $Q_{142}$, $Q_{12}$ и молярную теплоёмкость $C_{12}$. Выразить через $T_1$.}
\textbf{1. Температуры} из $pV=RT$: $T_3=T_4=2T_1$, $T_2=4T_1$. Для $i=5$: $C_V=\frac52R$, $C_p=\frac72R$.

\textbf{2. Путь 1--3--2:} изохора 1--3, $Q=C_V(2T_1-T_1)=2{,}5RT_1$; изобара 3--2, $Q=C_p(4T_1-2T_1)=7RT_1$. Сумма $9{,}5RT_1$.

\textbf{3. Путь 1--4--2:} изобара 1--4, $Q=C_p\cdot T_1=3{,}5RT_1$; изохора 4--2, $Q=C_V\cdot2T_1=5RT_1$. Сумма $8{,}5RT_1$.

\textbf{4. Путь 1--2:} работа — площадь трапеции под отрезком: $A=\frac{p_1+2p_1}2V_1=1{,}5RT_1$; $\Delta U=C_V(4T_1-T_1)=7{,}5RT_1$; $Q=\Delta U+A=9RT_1$.

\textbf{5. Теплоёмкость на 1--2.} Прямая через начало координат: $p=\alpha V$. Тогда $RT=pV=\alpha V^2$, дифференцируем: $R\,\dd T=2\alpha V\dd V=2p\,\dd V$,
т.е.\ $p\frac{\dd V}{\dd T}=\frac R2$ и $C=C_V+p\frac{\dd V}{\dd T}=\frac52R+\frac12R$.
\answer{$Q_{132}=9{,}5RT_1$, $Q_{142}=8{,}5RT_1$, $Q_{12}=9RT_1$, $C_{12}=3R$.}
\end{problem}

%-------------------------------------------------------------
\begin{problem}{1.5}{exdz}{Политропным называют обратимый процесс при постоянной теплоёмкости $C$. Найти уравнение политропы для моля идеального газа с известной $C_V$.}
\textbf{1. Теплоёмкость из первого начала:} $C\,\dd T=C_V\dd T+p\,\dd V$, $p=\frac{RT}V$:
\[(C-C_V)\,\dd T=\frac{RT}{V}\dd V\ \Rightarrow\ (C-C_V)\frac{\dd T}{T}=R\frac{\dd V}{V}.\]

\textbf{2. Интегрируем} ($C=$\,const): $(C-C_V)\ln T=R\ln V+\text{const}\Rightarrow T^{C_V-C}V^R=\text{const}$.

\textbf{3. Переходим к $p,V$:} возводим в степень $\frac1{C_V-C}$: $TV^{R/(C_V-C)}=\text{const}$; подставляем $T=\frac{pV}R$:
$pV^{n}=\text{const}$, где $n=1+\frac{R}{C_V-C}=\frac{C_V+R-C}{C_V-C}=\frac{C-C_p}{C-C_V}$.
\answer{$pV^n=\text{const}$ (и $TV^{n-1}=\text{const}$), $n=\dfrac{C-C_p}{C-C_V}$; обратно $C=C_V\dfrac{n-\gamma}{n-1}$.}
\note{$n=0$ — изобара ($C_p$), $n=1$ — изотерма ($C=\infty$), $n=\gamma$ — адиабата ($C=0$), $n=\infty$ — изохора ($C_V$); при $1<n<\gamma$ теплоёмкость $C<0$.}
\end{problem}

%-------------------------------------------------------------
\begin{problem}{1.6}{exdz}{Боковые стенки цилиндра, его крышка и невесомый поршень не проводят тепло, дно проводит. Над и под поршнем — по молю одного и того же
идеального газа с известной $C_V$. Газ в нижней части нагревают, поршень перемещается. Найти теплоёмкости $C_1$ (нижний газ) и $C_2$ (верхний) через их объёмы $V_1$, $V_2$.}
\textbf{1. Верхний газ} теплоизолирован, процесс медленный — адиабата: $C_2=0$ и $pV_2^\gamma=\text{const}$.
Дифференцируем логарифм: $\frac{\dd p}{p}=-\gamma\frac{\dd V_2}{V_2}$. Объём цилиндра постоянен: $\dd V_2=-\dd V_1$, значит $\frac{\dd p}{p}=\gamma\frac{\dd V_1}{V_2}$.

\textbf{2. Нижний газ.} Давления по обе стороны невесомого поршня равны. Уравнение состояния $pV_1=RT_1$, дифференцируем:
\[p\,\dd V_1+V_1\dd p=R\,\dd T_1\ \Rightarrow\ p\,\dd V_1\Big(1+\gamma\frac{V_1}{V_2}\Big)=R\,\dd T_1\ \Rightarrow\ p\frac{\dd V_1}{\dd T_1}=\frac{RV_2}{V_2+\gamma V_1}.\]

\textbf{3. Теплоёмкость} $C_1=C_V+p\frac{\dd V_1}{\dd T_1}$.
\answer{$C_1=C_V+R\dfrac{V_2}{V_2+\gamma V_1}$,\; $C_2=0$.}
\end{problem}

%-------------------------------------------------------------
\begin{problem}{1.7}{ex}{Найти разность теплоёмкостей $C_p-C_V$ для любого однородного изотропного вещества с известными уравнением состояния $p(V,T)$ и внутренней энергией $U(T,V)$.}
\textbf{1. Первое начало} с полным дифференциалом $U$:
$\delta Q=\dd U+p\,\dd V=\big(\frac{\partial U}{\partial T}\big)_V\dd T+\big[\big(\frac{\partial U}{\partial V}\big)_T+p\big]\dd V$.

\textbf{2.} Делим на $\dd T$. При $V=$\,const второе слагаемое пропадает: $C_V=\big(\frac{\partial U}{\partial T}\big)_V$.
При $p=$\,const: $C_p=C_V+\big[\big(\frac{\partial U}{\partial V}\big)_T+p\big]\big(\frac{\partial V}{\partial T}\big)_p$.
\answer{$C_p-C_V=\Big[\Big(\dfrac{\partial U}{\partial V}\Big)_T+p\Big]\Big(\dfrac{\partial V}{\partial T}\Big)_p$; для идеального газа ($\frac{\partial U}{\partial V}=0$, $p\frac{\partial V}{\partial T}=R$) это $R$.}
\end{problem}

%-------------------------------------------------------------
\begin{problem}{1.8}{exdz}{Обратимый цикл с молем идеального газа (рис.~1.9): 1--2 — изобара ($T_1\to T_2$), 2--3 — изохора ($T_2\to T_3$), 3--1 — адиабата.
Найти работу и теплоту на каждом этапе и КПД цикла как функцию $T_1$, $T_2$, $T_3$. $\gamma$ известна.}
\textbf{1. По участкам} (первое начало $Q=\Delta U+A$, $\Delta U=C_V\Delta T$):
\begin{itemize}
\item 1--2, изобара: $A_{12}=p_1(V_2-V_1)=R(T_2-T_1)$, $Q_{12}=C_p(T_2-T_1)>0$ — тепло получено;
\item 2--3, изохора: $A_{23}=0$, $Q_{23}=C_V(T_3-T_2)<0$ — тепло отдано;
\item 3--1, адиабата: $Q_{31}=0$, $A_{31}=-\Delta U=C_V(T_3-T_1)$.
\end{itemize}

\textbf{2. КПД:} $A=Q_{12}+Q_{23}$, тепло получено только на 1--2:
\[\eta=\frac{A}{Q_{12}}=1-\frac{C_V(T_2-T_3)}{C_p(T_2-T_1)}.\]
\answer{$\eta=1-\dfrac{T_2-T_3}{\gamma(T_2-T_1)}$.}
\end{problem}

%-------------------------------------------------------------
\begin{problem}{1.9}{ex}{Найти КПД цикла Карно (две изотермы $T_1>T_2$ и две адиабаты) с идеальным газом.}
\textbf{1. Теплоты на изотермах} ($\Delta U=0$, $Q=A$): $Q_{12}=\nu RT_1\ln\frac{V_2}{V_1}>0$, $Q_{34}=\nu RT_2\ln\frac{V_4}{V_3}<0$.

\textbf{2. Связь объёмов} — из адиабат 2--3 и 4--1 ($TV^{\gamma-1}=$\,const): $T_1V_2^{\gamma-1}=T_2V_3^{\gamma-1}$, $T_1V_1^{\gamma-1}=T_2V_4^{\gamma-1}$.
Делим одно на другое: $\frac{V_2}{V_1}=\frac{V_3}{V_4}$.

\textbf{3. КПД:} $\eta=1+\frac{Q_{34}}{Q_{12}}=1-\frac{T_2\ln(V_3/V_4)}{T_1\ln(V_2/V_1)}$, логарифмы сокращаются.
\answer{$\eta=1-\dfrac{T_2}{T_1}$.}
\end{problem}

%-------------------------------------------------------------
\begin{problem}{1.10}{ex}{Машина Карно с КПД 40\% используется как холодильная при тех же резервуарах. Какое количество теплоты $Q_2$ она переносит от холодильника
к нагревателю за цикл, если за цикл совершается работа $A=10$\,кДж? Оценить её эффективность.}
\textbf{1.} При обращении цикла все теплоты меняют знак, но по модулю те же. Для тепловой машины: $\eta=\frac A{Q_1}$, $A=Q_1-|Q_2|$.

\textbf{2.} Исключаем $Q_1=\frac A\eta$: $|Q_2|=Q_1-A=A\big(\frac1\eta-1\big)=10\cdot(2{,}5-1)=15$\,кДж.

\textbf{3. Холодильный коэффициент} $\eta_x=\frac{|Q_2|}{A}$.
\answer{$|Q_2|=15$ кДж,\; $\eta_x=1{,}5$.}
\end{problem}

%-------------------------------------------------------------
\begin{problem}{1.25}{dz}{Найти уравнение процесса для идеального газа, при котором молярная теплоёмкость зависит от температуры как $C=\alpha T$, $\alpha$ — постоянная. $C_V$ известна.}
\textbf{1. Теплоёмкость в процессе} (как в 1.5): $C=C_V+p\frac{\dd V}{\dd T}=C_V+\frac{RT}{V}\frac{\dd V}{\dd T}$. Приравниваем $\alpha T$:
\[\alpha T-C_V=\frac{RT}{V}\frac{\dd V}{\dd T}\ \Rightarrow\ R\frac{\dd V}{V}=\alpha\,\dd T-C_V\frac{\dd T}{T}.\]

\textbf{2. Интегрируем:} $R\ln V=\alpha T-C_V\ln T+\text{const}$. Делим на $R$ и потенцируем: $V\cdot T^{C_V/R}\cdot e^{-\alpha T/R}=\text{const}$.
\answer{$VT^{C_V/R}\exp(-\alpha T/R)=\text{const}$.}
\end{problem}

%-------------------------------------------------------------
\begin{problem}{1.28}{dz}{Моль идеального газа нагревают в цилиндре под поршнем, удерживаемым пружиной (закон Гука). Стенки и поршень теплоизолированы, дно проводит тепло.
Начальный объём $V_0$ (пружина не деформирована) подобран так, что $p_0S^2=kV_0$, где $p_0$ — наружное атмосферное давление, $S$ — площадь поршня,
$k$ — жёсткость пружины. Найти теплоёмкость газа. Массой поршня пренебречь.}
\textbf{1. Уравнение процесса.} Поршень сместился на $x=\frac{V-V_0}S$, пружина давит с силой $kx$. Равновесие поршня:
\[pS=p_0S+kx\ \Rightarrow\ p=p_0+\frac{k}{S^2}(V-V_0)=p_0+\frac{p_0}{V_0}(V-V_0)=\frac{p_0}{V_0}V\]
(использовали $\frac k{S^2}=\frac{p_0}{V_0}$). Процесс $p=\alpha V$ — прямая через начало координат.

\textbf{2. Теплоёмкость.} Как в 1.4: $RT=\alpha V^2\Rightarrow R\,\dd T=2p\,\dd V\Rightarrow p\frac{\dd V}{\dd T}=\frac R2$.
\answer{$C=C_V+\dfrac R2$.}
\end{problem}

%-------------------------------------------------------------
\begin{problem}{1.33}{dz}{Обратимый цикл с молем идеального газа: две изотермы 1--2 ($T_1$) и 3--4 ($T_2$) и две политропы 2--3, 4--1 с теплоёмкостью $C_0$ (рис.~1.15).
Найти КПД, если при изотермическом расширении объём увеличивается вдвое.}
\textbf{1. Отношение объёмов на изотермах.} На политропе $T^{C_V-C_0}V^R=$\,const (см.\ 1.5), т.е.\ $V\propto T^{(C_0-C_V)/R}$.
Обе политропы идут между $T_1$ и $T_2$, поэтому $\frac{V_3}{V_2}=\frac{V_4}{V_1}$, откуда $\frac{V_3}{V_4}=\frac{V_2}{V_1}=2$.

\textbf{2. Теплоты по участкам:}
\begin{itemize}
\item 1--2 (расширение при $T_1$): $Q=RT_1\ln2>0$;
\item 2--3 (политропа $T_1\to T_2$): $Q=C_0(T_2-T_1)<0$;
\item 3--4 (сжатие при $T_2$): $Q=-RT_2\ln2<0$;
\item 4--1 (политропа $T_2\to T_1$): $Q=C_0(T_1-T_2)>0$ (при $C_0>0$).
\end{itemize}

\textbf{3. КПД.} Работа $A=\sum Q=R(T_1-T_2)\ln2$ (теплоты на политропах сокращаются). Получено тепла $Q_+=RT_1\ln2+C_0(T_1-T_2)$.
\answer{$\eta=\dfrac{R(T_1-T_2)\ln2}{C_0(T_1-T_2)+RT_1\ln2}$.}
\note{1.32 (вместо политроп — изобары) — то же с $C_0\to C_p$.}
\end{problem}

%-------------------------------------------------------------
\begin{problem}{1.34}{dz}{Определить, каким путём теоретически эффективнее повысить КПД машины Карно: увеличить температуру нагревателя $T_1$ на $\Delta T$ при
фиксированной $T_2$ или понизить температуру холодильника $T_2$ на такую же $\Delta T$ при фиксированной $T_1$.}
\textbf{1. Два варианта:} $\eta_a=1-\dfrac{T_2}{T_1+\Delta T}$ (греем нагреватель), $\eta_b=1-\dfrac{T_2-\Delta T}{T_1}$ (охлаждаем холодильник).

\textbf{2. Сравниваем} дроби: $\eta_b>\eta_a\iff\frac{T_2-\Delta T}{T_1}<\frac{T_2}{T_1+\Delta T}\iff(T_2-\Delta T)(T_1+\Delta T)<T_1T_2$.
Раскрываем скобки: $T_2\Delta T-T_1\Delta T-\Delta T^2<0$, т.е.\ $-\Delta T(T_1-T_2+\Delta T)<0$ — верно, так как $T_1>T_2$.
\answer{эффективнее понижать температуру холодильника.}
\end{problem}

%-------------------------------------------------------------
\begin{problem}{1.35}{dz}{Моль одноатомного идеального газа ($\gamma=5/3$) совершает цикл Карно между резервуарами $t_1=127^\circ$C и $t_2=27^\circ$C.
Наименьший объём газа в цикле $V_1=5$\,л, наибольший $V_2=20$\,л. Найти работу за цикл $A$, теплоту $Q_1$ от нагревателя и теплоту $Q_2$, отданную холодильнику.}
\textbf{1. Где какие объёмы.} Наименьший объём — начало изотермического расширения при $T_1=400$\,К. Наибольший — конец расширения по адиабате, т.е.\ уже при $T_2=300$\,К.
Объём $V_a$ в конце изотермы $T_1$ находим по адиабате $V_a\to V_2$:
\[T_1V_a^{\gamma-1}=T_2V_2^{\gamma-1}\ \Rightarrow\ V_a=V_2\Big(\frac{T_2}{T_1}\Big)^{1/(\gamma-1)}=20\cdot0{,}75^{3/2}\approx13\text{ л}.\]

\textbf{2. Тепло от нагревателя} (изотерма $T_1$, $V_1\to V_a$):
$Q_1=RT_1\ln\frac{V_a}{V_1}=8{,}31\cdot400\cdot\ln2{,}6\approx3{,}17$\,кДж.

\textbf{3. Карно:} $\frac{|Q_2|}{Q_1}=\frac{T_2}{T_1}\Rightarrow Q_2=-\frac{300}{400}Q_1=-2{,}38$\,кДж; $A=Q_1+Q_2$.
\answer{$Q_1\approx3{,}17$ кДж, $Q_2\approx-2{,}38$ кДж, $A\approx0{,}79$ кДж.}
\end{problem}

%-------------------------------------------------------------
\begin{problem}{1.36}{dz}{Определить КПД цикла 1--2--3--4--1, в котором одноатомный идеальный газ последовательно проходит состояния $1(2p,V)$, $2(2p,8V)$, $3(p,4V)$, $4(p,2V)$.
Все участки цикла — политропы.}
\textbf{1. Температуры} в единицах $\frac{pV}{\nu R}$ (произведение координат): $T_1=2$, $T_2=16$, $T_3=4$, $T_4=2$.

\textbf{2. Тип каждого участка} (из $p_aV_a^n=p_bV_b^n$) и теплота $Q=\nu C\Delta T$, для одноатомного $C_V=\frac32R$, $C_p=\frac52R$:
\begin{itemize}
\item 1--2: $p=2p$ — изобара; $Q=\frac52\cdot(16-2)=35$;
\item 2--3: $\frac{p}{V}$ одинаково ($\frac{2p}{8V}=\frac{p}{4V}$) — процесс $p=\alpha V$, $C=C_V+\frac R2=2R$; $Q=2\cdot(4-16)=-24$;
\item 3--4: $p$ — изобара; $Q=\frac52\cdot(2-4)=-5$;
\item 4--1: $pV=2pV$ одинаково — изотерма, сжатие $2V\to V$; $Q=RT\ln\frac{V}{2V}=2\ln\frac12=-2\ln2$.
\end{itemize}
(все $Q$ в единицах $pV$.)

\textbf{3. КПД.} $A=\sum Q=35-24-5-2\ln2=6-2\ln2$; тепло получено только на 1--2: $Q_+=35$.
\answer{$\eta=\dfrac{6-2\ln2}{35}\approx13\%$.}
\end{problem}

%-------------------------------------------------------------
\begin{problem}{1.37}{dz}{Найти КПД обратимого цикла Отто из адиабат 1--2, 3--4 и изохор 2--3, 4--1 (рис.~1.16), рабочее тело — идеальный газ. Выразить через $T_1$ и $T_2$.}
\textbf{1. Теплоты.} На адиабатах $Q=0$. Изохора 4--1 — нагрев, тепло получено: $Q_1=C_V(T_1-T_4)$. Изохора 2--3 — охлаждение: $|Q_2|=C_V(T_2-T_3)$.

\textbf{2. Связь температур.} Обе адиабаты идут между одними и теми же объёмами $V_{\min}$ и $V_{\max}$ ($TV^{\gamma-1}=$\,const):
$\frac{T_2}{T_1}=\frac{T_3}{T_4}=\big(\frac{V_{\min}}{V_{\max}}\big)^{\gamma-1}$. Значит $T_2=kT_1$, $T_3=kT_4$ и $\frac{T_2-T_3}{T_1-T_4}=k=\frac{T_2}{T_1}$.

\textbf{3.} $\eta=1-\frac{|Q_2|}{Q_1}=1-\frac{T_2-T_3}{T_1-T_4}$.
\answer{$\eta=1-\dfrac{T_2}{T_1}$.}
\end{problem}
